% ======================================================================
%  Section insert: Eisenhart's lifts in detail  (+ comparison table)
%  Drop into hertz_chapter.tex in place of the old "Three geometrizations
%  side by side" section (i.e. before \section{Assessment...}).
%
%  Needs: amsmath, amsthm (definition/theorem/proposition/remark/example
%  environments as defined in the chapter), enumitem, booktabs.
%  Cross-references used: thm:hz-cyclic, sec:hz-hidden, sec:hz-jacobi,
%  prop:hz-dp-metric, eq:hz-V, eq:hz-KM, eq:hz-warped.
%  Bibliography keys used: hdp:eisenhart1929, hdp:duval1985,
%  hdp:duval1991, hdp:casetti2000.
% ======================================================================

\section{Eisenhart's lifts in detail}
\label{sec:hz-eisenhart}

In 1928--29 L.~P.~Eisenhart showed that the trajectories of a natural
Lagrangian system can be obtained as geodesics of a metric on a larger space
\cite{hdp:eisenhart1929}. His paper contains \emph{two} constructions,
and they stand in different relations to Hertz. The first is
Lorentzian, works for every potential (including time-dependent ones) and
needs two extra coordinates; the second is Riemannian, needs a positive
potential and one extra coordinate. The second turns out to be
the metric of Theorem~\ref{thm:hz-cyclic}. The first was rediscovered
in 1984--85 by Duval, Burdet, K\"unzle and Perrin, who called the resulting
spaces \emph{Bargmann structures} \cite{hdp:duval1985,hdp:duval1991}, and
the construction is now often called the Eisenhart--Duval lift.

\subsection{The Lorentzian (null) lift}

\begin{definition}[Eisenhart lift]\label{def:hz-eis}
Let $(Q,M)$ be a Riemannian manifold and $V\in C^\infty(Q)$ arbitrary. On
$\widehat Q=Q\times\Rb_t\times\Rb_z$ put
\begin{equation}\label{eq:hz-eis-metric}
  G_{\rm E}=M_{ab}(q)\,\ud q^a\ud q^b+2\,\ud t\,\ud z-2V(q)\,\ud t^2 .
\end{equation}
\end{definition}

The $(t,z)$ block $\bigl(\begin{smallmatrix}-2V&1\\1&0\end{smallmatrix}\bigr)$
has determinant $-1$ whatever $V$ is, so $G_{\rm E}$ is a non-degenerate metric
of Lorentzian signature $(n{+}1,1)$. Its inverse has the components
\begin{equation}\label{eq:hz-eis-inverse}
  G^{ab}=M^{ab},\qquad G^{tz}=1,\qquad G^{zz}=2V,\qquad G^{tt}=G^{ta}=G^{za}=0 .
\end{equation}
The only Christoffel symbols that differ from those of $M$ and from zero are
\begin{equation}\label{eq:hz-eis-christoffel}
  \Gamma^a_{tt}=M^{ab}\partial_bV,\qquad
  \Gamma^z_{at}=\Gamma^z_{ta}=-\partial_aV .
\end{equation}
(For the double pendulum this was checked symbolically for a flat base and
confirmed for the curved base $M$ through the Ricci tensor below.)

\begin{theorem}[Eisenhart]\label{thm:hz-eis}
A curve $\gamma(s)=(q(s),t(s),z(s))$ is a geodesic of $G_{\rm E}$ if and only if
\begin{equation}\label{eq:hz-eis-geod}
  \ddot t=0,\qquad
  \ddot q^a+\Gamma^a_{bc}\dot q^b\dot q^c=-M^{ab}\partial_bV\,\dot t^{\,2},\qquad
  \ddot z=2\,\dot t\,\partial_aV\,\dot q^a .
\end{equation}
Consequently, if $\dot t=1$ (so $s=t$ is physical time):
\begin{enumerate}[label=(\alph*),nosep]
  \item $q(t)$ solves Newton's equations $\ddot q^a+\Gamma^a_{bc}\dot q^b\dot q^c=-M^{ab}\partial_bV$;
  \item $\dot z=2V(q)-\varepsilon$ with a constant $\varepsilon$, so $p_t:=G(\partial_t,\dot\gamma)=-\varepsilon$;
  \item $G_{\rm E}(\dot\gamma,\dot\gamma)=2(E-\varepsilon)$, where $E=\tfrac12M\dot q\dot q+V$ is the energy of $q(t)$;
  \item $\gamma$ is a \emph{null} geodesic exactly when $\varepsilon=E$, and then
        $z(t)=z_0-\int_0^t L\,\ud t'=z_0-S[q]$ with $L=\tfrac12M\dot q\dot q-V$.
\end{enumerate}
Conversely every solution $q(t)$ of Newton's equations lifts, for every
choice of $\varepsilon$, $z_0$ and $t_0$, to a geodesic with $\dot t=1$.
\end{theorem}

\begin{proof}
Equations~\eqref{eq:hz-eis-geod} are the geodesic equations with the symbols
\eqref{eq:hz-eis-christoffel}: $\Gamma^t_{AB}=0$ gives $\ddot t=0$;
$\ddot q^a+\Gamma^a_{bc}\dot q^b\dot q^c+\Gamma^a_{tt}\dot t^2=0$ is the second;
$\ddot z+2\Gamma^z_{at}\dot q^a\dot t=0$ is the third. For $\dot t=1$ the third says
$\ddot z=2\,\ud V/\ud t$, hence $\dot z=2V-\varepsilon$, and $p_t=-2V\dot t+\dot z=-\varepsilon$.
Then
\[
  G(\dot\gamma,\dot\gamma)=M\dot q\dot q+2\dot z-2V
  =2T+2(2V-\varepsilon)-2V=2(T+V-\varepsilon)=2(E-\varepsilon),
\]
which gives (c) and the first half of (d). If $\varepsilon=E$ then
$\dot z=2V-E=V-T=-L$. The converse is immediate by defining
$z$ through $\dot z=2V-\varepsilon$.
\end{proof}

\begin{remark}[The extra coordinate is the action]\label{rem:hz-z-action}
Statement (d) is the geometric content of the lift: along the physical
(null) geodesics the new coordinate $z$ is minus Hamilton's principal
function evaluated along the path. This is no accident. With $\dot t=1$ the
geodesic Lagrangian is $\tfrac12G(\dot\gamma,\dot\gamma)=L+\dot z$, so
$\int\tfrac12G(\dot\gamma,\dot\gamma)\,\ud t=S[q]+z(t_2)-z(t_1)$:
the energy functional of the lift is Hamilton's action up to endpoint terms.
\end{remark}

\subsection{Hamiltonian form and conserved quantities}

The geodesic Hamiltonian $H_{\rm E}=\tfrac12G^{AB}p_Ap_B$ is, by~\eqref{eq:hz-eis-inverse},
\begin{equation}\label{eq:hz-eis-ham}
  H_{\rm E}=\tfrac12M^{ab}p_ap_b+p_tp_z+V(q)\,p_z^{\,2}.
\end{equation}
Both $t$ and $z$ are cyclic for time-independent $V$, so $p_t$ and $p_z$ are conserved.
With $p_z=G(\partial_z,\dot\gamma)=\dot t=1$ and
$H(q,p)=\tfrac12M^{ab}p_ap_b+V$ the physical Hamiltonian, \eqref{eq:hz-eis-ham} becomes
\[
  H_{\rm E}=H(q,p)+p_t ,
\]
and the null condition $H_{\rm E}=0$ is $p_t=-H=-E$: Hamilton's equations of $H_{\rm E}$
are the equations of the \emph{extended phase space}, in which energy is the momentum conjugate
to time and the dynamics is a constraint. The constant $p_z$ plays the role of a mass
(in the Bargmann picture it is the mass or the electric charge of the particle), and the
fixed value $p_z=1$ is what selects Newtonian dynamics with unit
inertia. Hamilton's equations give $\dot t=\partial H_{\rm E}/\partial p_t=p_z=1$ and
$\dot z=\partial H_{\rm E}/\partial p_z=p_t+2V=2V-E$, as found above.

\subsection{Geometry of the lift: a pp-wave}

\begin{proposition}[Curvature of the Eisenhart metric]\label{prop:hz-eis-curv}
\leavevmode
\begin{enumerate}[label=(\roman*),nosep]
  \item $\nabla\,\partial_z=0$: the null vector field $\partial_z$ is covariantly constant.
  \item The Ricci tensor of $G_{\rm E}$ has only the components
        \[
          \operatorname{Ric}_{ab}=\operatorname{Ric}^{M}_{ab},\qquad
          \operatorname{Ric}_{tt}=\Lap_MV .
        \]
\end{enumerate}
\end{proposition}

\begin{proof}
(i) is $\Gamma^A_{Bz}=0$, read off from \eqref{eq:hz-eis-christoffel}. For (ii) I computed the Ricci
tensor symbolically for a flat base with arbitrary $V(x,y)$ and, for the
unit-parameter double pendulum with its curved metric $M$, at random points, where the
only non-vanishing components were $\operatorname{Ric}_{ab}=K_MM_{ab}$ and
$\operatorname{Ric}_{tt}=\Lap_MV$ to rounding error; the general statement is the standard curvature
property of Eisenhart metrics \cite{hdp:casetti2000}.
\end{proof}

A Lorentzian manifold with a covariantly constant null vector whose metric has this
block form is a (generalized) \emph{pp-wave}, and the structure
$(\widehat Q,G_{\rm E},\partial_z)$ is a Bargmann structure \cite{hdp:duval1985,hdp:duval1991}.
Two consequences are worth stating.
\begin{itemize}[nosep]
  \item \emph{Gravity as curvature.} With $\Lap_MV=4\pi G\rho$ for a Newtonian
        potential, $\operatorname{Ric}_{tt}=\Lap_MV$ is the Poisson equation, and Ricci-flatness of the lift
        (flat $M$, harmonic $V$) is Newtonian gravity in vacuum. The lift makes Newtonian
        gravity \emph{look} like general relativity; it is a statement about a Lorentzian
        space of one dimension more than spacetime, not about spacetime.
  \item \emph{Tidal tensor.} The geodesic deviation equation of $G_{\rm E}$ along a lifted
        trajectory, restricted to the $q$-directions, is the usual variational equation of the
        dynamics,
        \[
          \frac{D^2\xi^a}{\ud t^2}+R^{M\,a}{}_{bcd}\,\dot q^b\xi^c\dot q^d
          +M^{ab}\nabla_b\nabla_cV\,\xi^c=0,
        \]
        and the trace of the tidal tensor is $\Lap_MV$ (the origin of $\operatorname{Ric}_{tt}$)
        \cite{hdp:casetti2000}. In contrast to the Jacobi metric of
        Section~\ref{sec:hz-jacobi}, no factor $E-V$ enters and there is no Hill-boundary singularity;
        but the price is that the criterion for exponential instability is the standard
        linear-stability computation, not a purely geometric sign condition.
\end{itemize}

\subsection{Examples}

\begin{example}[Harmonic oscillator]\label{ex:hz-eis-ho}
For $Q=\Rb$, $M=\ud q^2$ and $V=\tfrac12\omega^2q^2$,
$G_{\rm E}=\ud q^2+2\,\ud t\,\ud z-\omega^2q^2\ud t^2$ is a plane wave with $\operatorname{Ric}_{tt}=\omega^2$.
The null geodesics with $\dot t=1$ are $q=A\cos\theta$, $\theta=\omega t+\phi$, and
\[
  z=z_0+\tfrac14A^2\omega\,\sin2\theta ,
\]
which is $z_0-\int L\,\ud t$ with $L=-\tfrac12A^2\omega^2\cos2\theta$. The $z$-motion is a
bounded oscillation at twice the frequency.
\end{example}

\begin{example}[Double pendulum]\label{ex:hz-eis-dp}
With $M$ from Proposition~\ref{prop:hz-dp-metric} and $V$ from~\eqref{eq:hz-V},
\[
\begin{aligned}
  G_{\rm E}={}&(m_1+m_2)\,l_1^2\,\ud\theta_1^2+2m_2l_1l_2\cos(\theta_1-\theta_2)\,\ud\theta_1\ud\theta_2
  +m_2l_2^2\,\ud\theta_2^2\\
  &+2\,\ud t\,\ud z+2g\bigl[(m_1+m_2)l_1\cos\theta_1+m_2l_2\cos\theta_2\bigr]\ud t^2 .
\end{aligned}
\]
This is a four-dimensional Lorentzian manifold $\Tor^2\times\Rb^2$ with
$\operatorname{Ric}_{ab}=K_MM_{ab}$ (with $K_M$ from~\eqref{eq:hz-KM}) and $\operatorname{Ric}_{tt}=\Lap_MV$.
Its Killing fields include $\partial_z$ (null, parallel) and $\partial_t$; the rotation
$\partial_{\theta_1}+\partial_{\theta_2}$ that was a symmetry of the free pendulum is
broken by $V$. The null geodesics with $\dot t=1$ project to the chaotic trajectories of the pendulum,
and $z$ along them is $-S[\theta]$.
\end{example}

\subsection{The Riemannian lift, and its identity with the hidden rotor}

Eisenhart's second construction, for \emph{time-independent potentials that are positive} on $Q$,
uses one extra coordinate and a Riemannian metric of the form
$M+A(q)\,\ud z^2$ with $A$ built from $V$
\cite{hdp:eisenhart1929,hdp:casetti2000}. Taking $A=1/V_c$ with $V_c=V+c>0$,
\begin{equation}\label{eq:hz-eis-riem}
  G_{\rm R}=M_{ab}\,\ud q^a\ud q^b+\frac{1}{V_c(q)}\,\ud z^2 ,
\end{equation}
the Hamiltonian is $\tfrac12M^{ab}p_ap_b+\tfrac12V_c\,p_z^2$, which equals
the physical Hamiltonian $\tfrac12M^{ab}p_ap_b+V_c$ for $p_z=\sqrt2$. This is exactly
the metric~\eqref{eq:hz-warped} of Theorem~\ref{thm:hz-cyclic} with $f^2=p^2/(2V_c)$, $p=\sqrt2$,
and $\varphi=z$. The theorem therefore \emph{is} Eisenhart's Riemannian lift, with a different proof and with
a physical interpretation of the extra coordinate added: Eisenhart's $z$ is Hertz's hidden
cyclic coordinate, and $V=\tfrac12f^2\dot z^2$ is the kinetic energy of the hidden rotor.
(Hertz did not write down the warped metric; the idea that a potential is the kinetic energy of
a cyclic hidden motion is his, and the exact geometric statement is Eisenhart's.)

\subsection{Hertzian or not?}

Does the Lorentzian lift satisfy Hertz's requirements? Only in part. Its dynamics is
free in the sense that the law is geodesic, and no force appears. But Hertz's mass metric is positive
definite: kinetic energy is $\tfrac12\sum m\dot x^2\ge0$, and the straightest path is defined by minimizing
a norm of acceleration. In signature $(n{+}1,1)$ there is no such norm, the physical
trajectories are \emph{null}, and the extra coordinate $z$ is a null direction, not a hidden mass.
The Lorentzian lift is thus forceless but not Hertzian; the Riemannian lift is both.
That the two lifts are available side by side underlines the point of
Section~\ref{sec:hz-hidden}(c): the geometric content of the visible dynamics does not fix its higher-dimensional
realization.

\section{Four geometrizations side by side}
\label{sec:hz-compare}

\begin{table}[t]
\centering
\small
\setlength{\tabcolsep}{4pt}
\begin{tabular}{@{}lccc@{}}
\toprule
 & Hertz / Eisenhart--Riem. & Jacobi & Eisenhart--Lorentz\\
\midrule
Signature & Riemannian & Riemannian & Lorentzian $(n{+}1,1)$\\
Dimension & $n+1$ & $n$ & $n+2$\\
Fixed data & $p_z$ (hidden momentum) & energy $E$ & $p_z=1$, null condition\\
Restriction on $V$ & $V+c>0$ & Hill region $V<E$ & none\\
Geodesic parameter & physical time $t$ & $\tau$, $\ud\tau=(E-V)\ud t$ & physical time $t$\\
Where $V$ sits & hidden kinetic energy & conformal factor & $G_{tt}$\\
Role of extra coordinates & hidden angle $z$ & none & $t$ and null $z=-S$\\
Time-dependent $V$ & no & no & yes\\
Hertzian reading & hidden masses & none & none (null)\\
\bottomrule
\end{tabular}
\caption{Three ways (four constructions, since the first column is Hertz's idea and Eisenhart's
Riemannian metric at once) to turn the potential force into geometry. Only the first has a
mechanical reading in terms of positive masses.}
\label{tab:hz-compare}
\end{table}

Table~\ref{tab:hz-compare} summarizes the comparison. The Lorentzian lift is the only one that
treats time-dependent potentials uniformly, since $t$ is a coordinate of the lifted
space; the Jacobi construction requires conservation of energy and the Riemannian lift
requires a time-independent positive $V$.

